% Requires: \usepackage{graphicx} \usepackage{cleveref}
% The longer version is kept in section_diversity_long.tex.

\subsection{Diversity of the MAEMM's samples}
\label{sec:verb-diversity}

We measure the variety of a feature's eight samples by their Vendi score: the effective number of
distinct texts, from $1$ (all alike) to $8$ (all different). The MAEMM gives $2.2$, fewer than the
feature's top corpus windows ($2.9$) or an LLM writing from them ($3.4$), and about as many as NLA-AV
($2.0$; \cref{fig:verb-diversity-embed}a). This is not templating: the MAEMM's samples are the most
specific to their feature of any source (\cref{app:verb-diversity}). It gives a sharp but narrow
reading of each feature.

Narrow is not wrong. Where its samples agree most, the MAEMM fails on $4\%$ of features; where they
agree least, on $19\%$ (\cref{fig:verb-diversity-embed}b). Nor is spread a goal in itself: LoRAs
trained further on rare features spread more ($4.2$--$4.4$) only because their samples drift off the
feature.

\begin{figure}[htbp]
\centering
\includegraphics[width=\linewidth]{appendix/eval-inversion/fig11_diversity_embed.pdf}
\caption{\textbf{(a)} Distinct texts among eight samples per feature (Vendi score over bge-small
embeddings): median and interquartile range over the 131k-SAE 2k set. NLA-AV covers $500$ features,
the LLM $159$. \textbf{(b)} Share of features the MAEMM fails to verbalize
($\texttt{norm\_act} < 0.10$, best-of-8), by quartile of agreement among its samples.}
\label{fig:verb-diversity-embed}
\end{figure}
